% Единственный источник содержания темы. UTF-8.
% Метки label и ключи в [...] сохраняют опубликованные ссылки.

\section{Теория}

Теория предыдущих занятий:

\previoustheory

\section{Задачи}

\begin{task}[problem-solving-practice/problems/1]\label{task-problems-1}
\textbf{1.} Вероятность того, что в случайный момент времени температура тела здорового человека окажется ниже чем $36{,}8\,^{\circ}\mathrm{C}$, равна $0{,}81$. Найдите вероятность того, что в случайный момент времени у здорового человека температура окажется $36{,}8\,^{\circ}\mathrm{C}$ или выше.
\end{task}

\begin{solution}[problem-solving-practice/problems/1]\label{solution-problem-solving-practice-problems-1}
Температура «ниже $36{,}8^\circ\mathrm C$» и температура «$36{,}8^\circ\mathrm C$ или выше» — противоположные события. Поэтому искомая вероятность равна

\[1-0{,}81=0{,}19.\]

\textbf{Ответ:} $0{,}19$.
\end{solution}

\begin{task}[problem-solving-practice/problems/2]\label{task-problems-2}
\textbf{2.} При изготовлении подшипников диаметром 67 мм вероятность того, что диаметр будет отличаться от заданного не больше чем на 0,01 мм, равна 0,965. Найдите вероятность того, что случайный подшипник будет иметь диаметр меньше чем 66,99 мм или больше чем 67,01 мм.
\end{task}

\begin{solution}[problem-solving-practice/problems/2]\label{solution-problem-solving-practice-problems-2}
Диаметр отличается от $67$ мм не больше чем на $0{,}01$ мм, когда он находится в отрезке $[66{,}99;67{,}01]$. Требуемое событие противоположно этому, поэтому

\[1-0{,}965=0{,}035.\]

\textbf{Ответ:} $0{,}035$.
\end{solution}

\begin{task}[problem-solving-practice/problems/3]\label{task-problems-3}
\textbf{3.} Вероятность того, что батарейка бракованная, равна $0{,}06$. Покупатель в магазине выбирает случайную упаковку, в которой две таких батарейки. Найдите вероятность того, что обе батарейки окажутся исправными.
\end{task}

\begin{solution}[problem-solving-practice/problems/3]\label{solution-problem-solving-practice-problems-3}
Вероятность исправности одной батарейки равна $1-0{,}06=0{,}94$. В принятой в задаче модели исправности батареек независимы. Вероятность того, что исправны обе:

\[0{,}94\cdot0{,}94=0{,}8836.\]

\textbf{Ответ:} $0{,}8836$.
\end{solution}

\begin{task}[problem-solving-practice/problems/4]\label{task-problems-4}
\textbf{4.} Если шахматист А. играет белыми фигурами, то он выигрывает у шахматиста Б. с вероятностью 0,52. Если А. играет чёрными, то А. выигрывает у Б. с вероятностью 0,3. Шахматисты А. и Б. играют две партии, причём во второй партии меняют цвет фигур. Найдите вероятность того, что А. выиграет оба раза.
\end{task}

\begin{solution}[problem-solving-practice/problems/4]\label{solution-problem-solving-practice-problems-4}
Одну партию А. играет белыми, другую — чёрными. В принятой в задаче модели результаты партий считаются независимыми. Поэтому вероятность двух побед равна

\[0{,}52\cdot0{,}3=0{,}156.\]

\textbf{Ответ:} $0{,}156$.
\end{solution}

\begin{task}[problem-solving-practice/problems/5]\label{task-problems-5}
\textbf{5.} Агрофирма закупает куриные яйца только в двух домашних хозяйствах. Известно, что 5\% яиц из первого хозяйства  —  яйца высшей категории, а из второго хозяйства  —  30\% яиц высшей категории. В этой агрофирме 15\% яиц высшей категории. Найдите вероятность того, что яйцо, купленное у этой агрофирмы, окажется из первого хозяйства.
\end{task}

\begin{solution}[problem-solving-practice/problems/5]\label{solution-problem-solving-practice-problems-5}
\textbf{Способ 1. Через условную вероятность}

Обозначим случайные события $A=\{\text{яйцо из первого хозяйства}\}$, $B=\overline A=\{\text{яйцо из второго хозяйства}\}$ и $C=\{\text{яйцо высшей категории}\}$. По условию

\[\mathbb P(C\mid A)=0{,}05,\qquad \mathbb P(C\mid\overline A)=0{,}3,\qquad \mathbb P(C)=0{,}15.\]

Пусть $\mathbb P(A)=x$. Тогда $\mathbb P(\overline A)=1-x$. Распишем определение условной вероятности:

\[\mathbb P(C\mid A)=\frac{\mathbb P(C\cap A)}{\mathbb P(A)}=\frac{\mathbb P(C\cap A)}{x}=0{,}05,\]

\[\mathbb P(C\mid\overline A)=\frac{\mathbb P(C\cap\overline A)}{\mathbb P(\overline A)}=\frac{\mathbb P(C\cap\overline A)}{1-x}=0{,}3.\]

Следовательно,

\[\mathbb P(C\cap A)=0{,}05x,\qquad \mathbb P(C\cap\overline A)=0{,}3(1-x).\]

События $C\cap A$ и $C\cap\overline A$ несовместны, а их объединение равно $C$: яйцо высшей категории поступило либо из первого, либо из второго хозяйства. Поэтому по формуле сложения вероятностей несовместных событий

\[\mathbb P(C)=\mathbb P(C\cap A)+\mathbb P(C\cap\overline A).\]

Подставляем известные значения:

\[0{,}15=0{,}05x+0{,}3(1-x)=0{,}3-0{,}25x.\]

Отсюда $0{,}25x=0{,}15$, поэтому $x=0{,}6$. Искомая вероятность $\mathbb P(A)=0{,}6$.

\textbf{Ответ:} $0{,}6$.

\textbf{Способ 2. Через количество яиц}

Пусть закуплено $N$ яиц, из них $x$ — из первого хозяйства. Тогда из второго хозяйства поступило $N-x$ яиц. Подсчитаем число яиц высшей категории:

\[0{,}05x+0{,}30(N-x)=0{,}15N.\]

Раскрывая скобки, получаем $0{,}25x=0{,}15N$, то есть $x=0{,}6N$. Значит, доля яиц из первого хозяйства и вероятность выбрать такое яйцо равны $0{,}6$. \textbf{Ответ:} $0{,}6$.
\end{solution}

\begin{task}[problem-solving-practice/problems/6]\label{task-problems-6}
\textbf{6.} Всем пациентам с подозрением на гепатит делают анализ крови. Если анализ выявляет гепатит, то результат анализа называется положительным. У больных гепатитом пациентов анализ даёт положительный результат с вероятностью $0,9$. Если пациент не болен гепатитом, то анализ может дать ложный положительный результат с вероятностью $0,01$. Известно, что $5\%$ пациентов, поступающих с подозрением на гепатит, действительно больны гепатитом. Найдите вероятность того, что результат анализа у пациента, поступившего в клинику с подозрением на гепатит, будет положительным.
\end{task}

\begin{solution}[problem-solving-practice/problems/6]\label{solution-problem-solving-practice-problems-6}
Обозначим через $H$ наличие гепатита, через $T$ — положительный результат анализа. События $H$ и $\overline H$ разбивают всех пациентов на две несовместные группы.

Вероятность положительного анализа у больного пациента:

\[\mathbb P(H\cap T)=0,05\cdot0,9=0,045.\]

Вероятность ложного положительного результата у здорового:

\[\mathbb P(\overline H\cap T)=0,95\cdot0,01=0,0095.\]

Складываем вероятности двух несовместных способов получить положительный результат:

\[\mathbb P(T)=0,045+0,0095=0,0545.\]
\end{solution}

\section{Домашнее задание}

\begin{task}[problem-solving-practice/homework/1]\label{task-homework-1}
\textbf{1.} Вероятность того, что новый электрический чайник прослужит больше года, равна 0,97. Вероятность того, что он прослужит больше двух лет, равна 0,89. Найдите вероятность того, что он прослужит меньше двух лет, но больше года.
\end{task}

\begin{solution}[problem-solving-practice/homework/1]\label{solution-problem-solving-practice-homework-1}
Обозначим срок службы чайника через $T$. Как принято в модели этой задачи, считаем, что вероятность поломки ровно через два года равна нулю: $\mathbb P(T=2)=0$.

Событие «прослужил больше года» состоит из несовместных случаев: $1<T<2$, $T=2$ и $T>2$. Поэтому

\[\mathbb P(1<T<2)=0{,}97-0-0{,}89=0{,}08.\]

\textbf{Ответ:} $0{,}08$.
\end{solution}

\begin{task}[problem-solving-practice/homework/2]\label{task-homework-2}
\textbf{2.} Из районного центра в деревню ежедневно ходит автобус. Вероятность того, что в понедельник в автобусе окажется меньше 18 пассажиров, равна 0,82. Вероятность того, что окажется меньше 10 пассажиров, равна 0,51. Найдите вероятность того, что число пассажиров будет от 10 до 17.
\end{task}

\begin{solution}[problem-solving-practice/homework/2]\label{solution-problem-solving-practice-homework-2}
Событие «меньше $18$ пассажиров» состоит из двух несовместных случаев: пассажиров меньше $10$ или их от $10$ до $17$ включительно. Следовательно, искомая вероятность равна

\[0{,}82-0{,}51=0{,}31.\]

\textbf{Ответ:} $0{,}31$.
\end{solution}

\begin{task}[problem-solving-practice/homework/3]\label{task-homework-3}
\textbf{3.} Помещение освещается фонарём с двумя лампами. Вероятность перегорания лампы в течение года равна 0,3. Найдите вероятность того, что в течение года хотя бы одна лампа не перегорит.
\end{task}

\begin{solution}[problem-solving-practice/homework/3]\label{solution-problem-solving-practice-homework-3}
Противоположное событие — обе лампы перегорят. В принятой в задаче модели перегорания ламп независимы, поэтому вероятность этого события равна $0{,}3^2=0{,}09$. Значит, хотя бы одна лампа не перегорит с вероятностью

\[1-0{,}09=0{,}91.\]

\textbf{Ответ:} $0{,}91$.
\end{solution}

\begin{task}[problem-solving-practice/homework/4]\label{task-homework-4}
\textbf{4.} Какова вероятность того, что случайно выбранный телефонный номер оканчивается двумя чётными цифрами?
\end{task}

\begin{solution}[problem-solving-practice/homework/4]\label{solution-problem-solving-practice-homework-4}
В обычной равновероятной модели последних двух цифр всего $10\cdot10=100$ пар: от $00$ до $99$. Чётных цифр пять: $0,2,4,6,8$. Обе цифры чётные в $5\cdot5=25$ парах, поэтому

\[\mathbb P=\frac{25}{100}=0{,}25.\]

\textbf{Ответ:} $0{,}25$.
\end{solution}

\begin{task}[problem-solving-practice/homework/5]\label{task-homework-5}
\textbf{5.} Ковбой Джон попадает в муху на стене с вероятностью 0,9, если стреляет из пристрелянного револьвера. Если Джон стреляет из непристрелянного револьвера, то он попадает в муху с вероятностью 0,2. На столе лежит 10 револьверов, из них только 4 пристрелянные. Ковбой Джон видит на стене муху, наудачу хватает первый попавшийся револьвер и стреляет в муху. Найдите вероятность того, что Джон промахнётся.
\end{task}

\begin{solution}[problem-solving-practice/homework/5]\label{solution-problem-solving-practice-homework-5}
Вероятность промаха из пристрелянного револьвера равна $1-0{,}9=0{,}1$, из непристрелянного — $1-0{,}2=0{,}8$.

Событие «выбран пристрелянный револьвер и произошёл промах» имеет вероятность $\frac4{10}\cdot0{,}1=0{,}04$. Событие «выбран непристрелянный револьвер и произошёл промах» имеет вероятность $\frac6{10}\cdot0{,}8=0{,}48$. Эти события несовместны и вместе составляют событие «промах», поэтому по формуле сложения

\[\mathbb P(\text{промах})=0{,}04+0{,}48=0{,}52.\]

\textbf{Ответ:} $0{,}52$.
\end{solution}

